\section{Domain-Specific Models (Workshops 2--5 Summary)}
\label{sec:domain_models}

Before analyzing cross-domain transferability, we summarize the exact architectures and persisted test results built across Workshops~2--5 on the \texttt{flwrlabs/shakespeare} domain.

\paragraph{Workshop 2: Autoregressive \texttt{CharTransformer} with Post-Hoc Logit Adjustment.}
Implemented in \texttt{shared/models.py}, our core sequence model is a 4-layer Pre-LN causal Transformer (\texttt{d\_model}=256, 4 attention heads, feedforward width 1,024, GELU activations, weight-tied input/output embeddings) trained on contiguous 81-character blocks with AdamW ($\beta_1=0.9, \beta_2=0.95$, weight decay 0.05), linear warmup + cosine annealing, and mixed precision. On the full 644,398-row unseen-character test split (\texttt{Workshop\_2/reports/results\_comparison.csv}), the calibrated Transformer ($\tau=0$) achieved \WTwoCalAcc{} top-1 accuracy, \WTwoCalTopFive{} top-5 accuracy, \WTwoCalBPC{}~BPC, \WTwoCalFOne{} macro-F1, \WTwoCalAUC{} macro OvR ROC-AUC, and \WTwoCalECE{} ECE—more than doubling W1 accuracy while halving cross-entropy. Applying post-hoc logit adjustment~\cite{Menon2020long} tuned strictly on validation ($\tau=0.10$) raised macro-recall from 29.30\% to 30.37\%, lifted macro-F1 to \WTwoAdjFOne{}, and improved calibration to \WTwoAdjECE{} ECE without retraining.

\paragraph{Workshop 3: Controllable Text Generation via Double+Dueling DQN.}
Workshop~3 wrapped the frozen Workshop~2 Transformer inside a Gymnasium environment (\texttt{TextGenerationEnv}, 290-dimensional observation combining 34 lexical candidate features with the L2-normalized 256-dimensional Transformer hidden state; 5 discrete candidate-character actions; 100-step episodes). A Double+Dueling Q-Network with Polyak soft updates ($\tau=0.005$) and Huber loss was trained over 400~episodes across four hyperparameter configurations. On 150 unseen test contexts across three seeds (\texttt{Workshop\_3/reports/test\_comparison.csv}), the DQN policy improved mean episodic return from \WThreeGreedyReturn{} (greedy Transformer) and $-72.811$ (random policy) to \textbf{\WThreeDQNReturn{}}, elevating novel valid word rate from 60.07\% to \WThreeDQNNewWord{} while maintaining a \WThreeDQNInvalid{} invalid-word rate and eliminating repetitive loops.

\paragraph{Workshop 4: 1D Residual Convolutional AE/VAE and Anomaly Detection.}
Workshop~4 implemented \texttt{SeqTextAutoencoder} using dilated 1D residual blocks (\texttt{GroupNorm} + GELU + \texttt{Conv1d} with dilations 1, 2, 4), learned decoder positional embeddings, and Feature-wise Linear Modulation (FiLM) from latent vector $z \in \mathbb{R}^{128}$. To prevent posterior collapse in the variational variant, training combined Free Bits ($0.25$ nats/dimension) with 4-cycle cyclical KL annealing and Monte Carlo ELBO checkpoint selection constrained to $\mathrm{KL} \ge 10.0$. On 20,000 test windows (\texttt{Workshop\_4/reports/reconstruction\_metrics.csv} and \texttt{anomaly\_detection\_metrics.csv}), the deterministic AE reached \WFourAEAcc{} character reconstruction accuracy ($\mathrm{CE}=0.653$) and the $\beta=0.5$ VAE reached \WFourVAEAcc{} ($\mathrm{CE}=1.095$). Both detected corrupted/out-of-distribution sequences at near-ceiling accuracy (VAE recon+KL PR-AUC $1.000$, F1 \WFourVAEAnomFOne{}), though linear probing (\texttt{silhouette\_scores.csv}) revealed that supervised Workshop~2 hidden states remained far more discriminative for next-character prediction (probe F1 0.511 vs.\ 0.196 for AE and 0.188 for VAE).

\paragraph{Workshop 5: Discrete 1D GAN / WGAN-GP and Controlled Data Augmentation.}
Workshop~5 constructed residual 1D Vanilla GAN (with spectral normalization and unilateral label smoothing $0.9$) and WGAN-GP (critic with \texttt{LayerNorm}, residual scaling $0.3h$, FP32 gradient penalty $\lambda_{\mathrm{GP}}=10$, 5~critic steps per generator step) using Gumbel-Softmax straight-through sampling ($\tau: 2.0 \to 0.5$) to backpropagate through discrete character outputs. Crucially, Workshop~5 produced an honest empirical finding (\texttt{Workshop\_5/reports/metrics/gan\_quality.csv} and \texttt{augmentation\_results.csv}): while both GANs avoided mode collapse (99.75\%--100\% unique sequences, zero 20-character memorization), discrete text fluency remained poor (LM perplexity \WFiveVanillaPPL{} for Vanilla GAN and \WFiveWGANPPL{} for WGAN-GP). Retraining the Workshop~2 Transformer with 10\% W2-filtered synthetic WGAN-GP blocks reduced test BPC marginally from \WFiveRealBPC{} (\texttt{real\_only}) to \WFiveWGANBPC{}, whereas simple domain-specific character noise augmentation achieved the strongest macro-F1 gain (\WFiveNoiseFOne{} vs.\ 0.3185). This lesson directly informs our cautious, evidence-driven design of Ablation~4 in Workshop~6.
